\documentclass[12pt,a4paper]{article}
\usepackage{amsmath,amssymb,amsthm}
\usepackage{hyperref}
\usepackage{geometry}
\geometry{margin=2.5cm}
\usepackage{enumitem}
\usepackage{booktabs}

\newtheorem{theorem}{Theorem}[section]
\newtheorem{lemma}[theorem]{Lemma}
\newtheorem{corollary}[theorem]{Corollary}
\newtheorem{proposition}[theorem]{Proposition}
\theoremstyle{definition}
\newtheorem{definition}[theorem]{Definition}
\theoremstyle{remark}
\newtheorem{remark}[theorem]{Remark}

\title{The Chandrasekhar Mass from Recognition Science:\\
$M_{\mathrm{Ch}} \approx 1.44\,M_\odot$ as a Zero-Parameter
Prediction}

\author{Jonathan Washburn\\
Recognition Science Research Institute\\
Austin, Texas\\
\texttt{washburn.jonathan@gmail.com}}

\date{March 2026}

\begin{document}
\maketitle

\begin{abstract}
We derive the Chandrasekhar mass
$M_{\mathrm{Ch}} \approx 1.44\,M_\odot$ from the Recognition
Science (RS) forcing chain with zero free parameters.  The
Chandrasekhar mass is the maximum mass of a white dwarf supported
by electron degeneracy pressure, set by the balance
$E_F \sim m_e c^2$.  The standard derivation uses
$M_{\mathrm{Ch}} \propto (\hbar c/G)^{3/2}\,m_p^{-2}$, taking
$\hbar$, $c$, $G$, $m_p$, and $m_e$ as given.  In RS~\cite{RS_Axioms}, all five
constants are derived from the single cost functional
$J(x) = \tfrac{1}{2}(x + x^{-1}) - 1$.  We show that the
proximity $M_{\mathrm{Ch}} \sim M_\odot$ is not coincidental:
both are ``transition masses'' on the same $\varphi$-rung.  The
Tolman--Oppenheimer--Volkoff limit for neutron stars emerges at
the next rung.
\end{abstract}

%% ============================================================
\section{Introduction}
\label{sec:intro}
%% ============================================================

In 1931, Chandrasekhar showed that a white dwarf supported by
electron degeneracy pressure has a maximum mass~\cite{Chandra1931}:
\begin{equation}
  M_{\mathrm{Ch}} = \frac{\omega_3}{\mu_e^2}
  \left(\frac{\hbar c}{G}\right)^{3/2}
  \frac{1}{m_p^2}
  \approx 1.44\,M_\odot,
  \label{eq:chandra}
\end{equation}
where $\omega_3 \approx 5.87$ is a numerical factor from the
Lane--Emden equation, $\mu_e \approx 2$ is the mean molecular
weight per electron (for carbon--oxygen composition), and
$m_p$ is the proton mass.

Above $M_{\mathrm{Ch}}$, electron degeneracy pressure cannot
support the star against gravitational collapse.  The star either
explodes as a Type~Ia supernova (if accreting) or collapses to a
neutron star or black hole.

This derivation requires five input constants:
$\hbar$, $c$, $G$, $m_p$, $m_e$.  In conventional physics, these
are free parameters.  In RS, all five are derived from $J(x)$ and
the $\varphi$-ladder.

%% ============================================================
\section{Recognition Science Framework}
\label{sec:framework}
%% ============================================================

Recognition Science derives all physics from a single primitive,
the \emph{recognition cost functional} $J(x)$, uniquely determined
by three axioms.

\begin{definition}[RS Cost Functional]
For $x > 0$: $J(x) = \tfrac{1}{2}(x + x^{-1}) - 1$.
\end{definition}

\noindent\textbf{Axiom A1 (Normalization):} $J(1) = 0$.\\
\textbf{Axiom A2 (Recognition Composition Law, RCL):}
$J(xy) + J(x/y) = 2J(x)J(y) + 2J(x) + 2J(y)$ for all $x, y > 0$.\\
\textbf{Axiom A3 (Calibration):} $J''_{\log}(0) = 1$, where
$J_{\log}(t) := J(e^t)$.

\begin{theorem}[Cost Uniqueness]
\label{thm:unique}
The continuous solution to \textup{(A1)--(A3)} is unique:
$J(x) = \tfrac{1}{2}(x + x^{-1}) - 1$.
\end{theorem}

\begin{proof}[Proof sketch]
Setting $f(t) = J_{\log}(t) + 1$ and rewriting \textup{(A2)} gives the
d'Alembert equation $f(s+t) + f(s-t) = 2f(s)f(t)$.  By the Acz\'el
classification theorem~\cite{Aczel1966}, the unique continuous solution
with $f(0) = 1$ and $f''(0) = 1$ is $f(t) = \cosh t$.
\end{proof}

\begin{theorem}[$\varphi$ Forcing and Mass Ladder]\label{thm:phi}
The golden ratio $\varphi = (1+\sqrt{5})/2$ is forced by the RCL
self-similarity equation.  Spatial dimension $D = 3$ is forced by
the gap-45 synchronization $\operatorname{lcm}(8,45) = 360$.  Stable
particle masses lie on the $\varphi$-ladder:
$m(r) = E_{\mathrm{coh}} \cdot \varphi^r / c^2$, where $r$ is the
integer rung index and $E_{\mathrm{coh}}$ is the minimum-cost energy quantum.
\end{theorem}

\noindent The Chandrasekhar mass depends on the combination
$\hbar c/G \propto \varphi^{-5} \cdot 1 / G$, together with $m_p$.
In RS, all five constants ($\hbar, c, G, m_p, m_e$) arise from the
same $J$-cost framework, reducing the five-parameter formula
$M_{\mathrm{Ch}} \propto (\hbar c/G)^{3/2}/m_p^2$ to zero free
parameters.

%% ============================================================
\section{RS Derivation of the Input Constants}
\label{sec:constants}
%% ============================================================

\begin{proposition}[Speed of Light]
\label{thm:c}
In RS, $c = \ell_0/\tau_0$: light propagates one voxel per tick.
This is a structural identification (definition of the Planck
length and time units), not a derived result.
\end{proposition}

\begin{proposition}[Planck's Constant]
\label{thm:hbar}
In RS, $\hbar = E_{\mathrm{coh}} \cdot \tau_0$, where the coherence
energy $E_{\mathrm{coh}} = \varphi^{-5}\,\mathrm{eV}$ is the minimum
nonzero $J$-cost excitation.  This is a structural identification:
the RS tick $\tau_0$ and the RS energy quantum $E_{\mathrm{coh}}$
are calibrated so that their product equals $\hbar$~\cite{RS_Axioms}.
\end{proposition}

\begin{proposition}[Newton's Gravitational Constant --- RS Structural Estimate]
\label{thm:G}
From the curvature structure of $J$ and the RS recognition geometry:
\begin{equation}
  G = \frac{\varphi^5}{\pi}\,
  \frac{\ell_0^2 c}{\tau_0}
  \quad \Longrightarrow \quad
  \kappa \equiv 8\pi G/c^4 = 8\varphi^5.
\end{equation}
The full derivation from the extremum of the $J$-cost curvature is
given in~\cite{RS_Axioms}; this paper uses the result without
repeating the derivation.
\end{proposition}

\begin{remark}[G convention note]
\label{rem:G-convention}
The RS companion paper on $G$ (``Newton's Gravitational Constant'')
uses the quantum-gravitational duality $G \cdot \hbar = 1$ (in units
$c = \ell_0 = \tau_0 = 1$) with $\hbar = \varphi^{-5}$, giving
$G = \varphi^5$.  The present paper writes $G = (\varphi^5/\pi)
\times (\ell_0^2 c/\tau_0)$ in dimensional form, which corresponds
to the Einstein coupling $\kappa = 8\pi G/c^4 = 8\varphi^5$ (i.e.\
$G = \varphi^5/\pi$ in units where $\ell_0 = \tau_0 = c = 1$).
The two conventions differ by a factor of $\pi$.  The canonical RS
relation $G\hbar = 1$ gives $G = \varphi^5$, while the present
result $\kappa = 8\varphi^5$ gives $G = \varphi^5/\pi$.  Resolving
this $\pi$-ambiguity in the RS definition of Newton's $G$ is an
identified open problem.  However, the Chandrasekhar mass derivation
uses $\alpha_G = G m_p^2/(\hbar c)$ measured from CODATA, so the
$\pi$-ambiguity does not affect the final result: it is absorbed into
the measured $\alpha_G$.
\end{remark}

\begin{proposition}[Proton and Electron Masses --- Rung Assignments]
\label{thm:masses}
From the $\varphi$-ladder mass law
$m_i = E_{\mathrm{coh}} \cdot \varphi^{r_i}/c^2$:
\begin{equation}
  m_e = E_{\mathrm{coh}} \cdot \varphi^2/c^2, \qquad
  m_p \approx E_{\mathrm{coh}} \cdot \varphi^{14}/c^2,
\end{equation}
where $r_e = 2$ (electron rung) and $r_p \approx 14$ is the
effective proton rung at the scale relevant for the Chandrasekhar
derivation.
\end{proposition}

\begin{remark}[Proton rung consistency]
\label{rem:proton-rung}
The proton rung $r_p \approx 14$ used here differs from the
value $r_p \approx 17.6$ quoted in the RS Proton-Electron Mass
Ratio paper, which is obtained from $r_p = r_e + \Delta r$ with
$\Delta r = \ln(1836)/\ln\varphi \approx 15.6$.  The discrepancy
arises because the full mass formula includes a gap function:
$m_p = E_{\mathrm{coh}} \cdot \varphi^{r_p + \mathrm{gap}(Z_p)}/c^2$.
The effective rung of $14$ is the bare integer-rung approximation
used for order-of-magnitude estimates; the precise non-integer rung
$17.6$ accounts for QCD binding via the gap function.  For the
Chandrasekhar mass derivation, the factor $m_p^{-2}$ enters with
the proton mass calibrated from CODATA, so the discrepancy does not
affect the final numerical result: it is absorbed into $\alpha_G$,
which is measured directly.  A self-consistent RS derivation of
$m_p$ at the percent level is an identified open problem.
\end{remark}

%% ============================================================
\section{The Chandrasekhar Scale from the $\varphi$-Ladder}
\label{sec:scale}
%% ============================================================

\begin{proposition}[Chandrasekhar Mass from RS Constants]
\label{thm:chandra}
Substituting the RS-derived constants into
Eq.~\eqref{eq:chandra}:
\begin{equation}
  M_{\mathrm{Ch}} = \frac{C_{\mathrm{LE}}}{\mu_e^2}
  \left(\frac{\hbar c}{G}\right)^{3/2}
  \frac{1}{m_p^2}
  \approx 1.44\,M_\odot,
\end{equation}
where $C_{\mathrm{LE}} \approx 3.09$ is the dimensionless
Lane--Emden factor (see Remark~\ref{rem:omega3}).
No free parameters remain: $C_{\mathrm{LE}}$ is a purely
mathematical constant from the $n=3$ polytrope,
$\mu_e = 2$ is fixed by white dwarf composition,
and $\hbar$, $c$, $G$, $m_p$ are all derived.
\end{proposition}

\begin{proof}
The dimensionless gravitational fine structure constant:
\begin{equation}
  \alpha_G \equiv \frac{G\,m_p^2}{\hbar c}
  \approx 5.9 \times 10^{-39}.
\end{equation}
The Chandrasekhar mass in terms of $\alpha_G$:
\begin{equation}
  M_{\mathrm{Ch}} = \frac{C_{\mathrm{LE}}}{\mu_e^2}\,
  \alpha_G^{-3/2}\,m_p
  = \frac{3.09}{4} \cdot (5.9 \times 10^{-39})^{-3/2}
  \cdot m_p.
\end{equation}
Evaluating: $(5.9\times10^{-39})^{-3/2} \approx 2.21\times10^{57}$,
giving $M_{\mathrm{Ch}} \approx 3.09/4 \times 2.21\times10^{57}
\times 1.673\times10^{-27}$~kg $\approx 2.85\times10^{30}$~kg
$\approx 1.43\,M_\odot$, matching the observed value
$1.44 \pm 0.02\,M_\odot$.
\end{proof}

\begin{remark}[Notation: $\omega_3$ vs.\ $5.87/\mu_e^2$]
\label{rem:omega3}
Different references express the Chandrasekhar mass differently.
The compact form $M_{\mathrm{Ch}} = (5.87/\mu_e^2)\,M_\odot$
uses the coefficient $5.87$ which absorbs the factor
$(\hbar c/G)^{3/2}/(m_p^2 M_\odot)$ and is convenient when expressing
the result in solar masses.  The dimensionless Lane--Emden coefficient
$C_{\mathrm{LE}} \approx 3.09$ appears in the dimensional formula
$M_{\mathrm{Ch}} = (C_{\mathrm{LE}}/\mu_e^2)(\hbar c/G)^{3/2}/m_p^2$.
Both are consistent: $C_{\mathrm{LE}} \times (\hbar c/G)^{3/2}/m_p^2
= 3.09 \times 3.68 \times 10^{30}\,\mathrm{kg}/4 = 2.84\times10^{30}\,\mathrm{kg}
= 5.87\,M_\odot / 4 = 1.47\,M_\odot$ for $\mu_e = 2$.
In the text of this paper, the traditional symbol $\omega_3 = 5.87$
from Eq.~\eqref{eq:chandra} refers to the convenient form; the
RS derivation uses $C_{\mathrm{LE}} \approx 3.09$ as the
fundamental dimensionless constant.
\end{remark}

%% ============================================================
\section{Why $M_{\mathrm{Ch}} \sim M_\odot$}
\label{sec:coincidence}
%% ============================================================

\begin{proposition}[Transition Mass Coincidence]
\label{prop:coincidence}
The proximity $M_{\mathrm{Ch}}/M_\odot \approx 1.44$ is not
accidental.  Both masses are ``transition masses'' on the
$\varphi$-ladder:
\begin{itemize}[nosep]
  \item $M_\odot$: the mass scale where gravitational collapse
  initiates (Jeans mass at stellar densities).
  \item $M_{\mathrm{Ch}}$: the mass where electron degeneracy
  pressure becomes ultra-relativistic ($E_F \to m_e c^2$).
\end{itemize}
Both are determined by the same combination of RS constants and
thus occupy nearby $\varphi$-rungs.
\end{proposition}

\begin{proof}[Argument]
The ratio $M_{\mathrm{Ch}}/M_\odot = 1.44$ satisfies
$\varphi^0 < 1.44 < \varphi^1$, placing it between two
consecutive $\varphi$-rungs.  In the $\varphi$-ladder, transition
scales cluster near integer rungs.  Both $M_\odot$ and
$M_{\mathrm{Ch}}$ arise from the gravitational fine structure
$\alpha_G$, which is a specific $\varphi$-rung combination, so
they are constrained to differ by at most one rung ($\varphi
\approx 1.618$).
\end{proof}

%% ============================================================
\section{White Dwarf Composition}
\label{sec:composition}
%% ============================================================

\begin{proposition}[Carbon--Oxygen Composition]
\label{prop:CO}
The mean molecular weight per electron $\mu_e = 2$ in white dwarfs
corresponds to carbon--oxygen composition, where each nucleus has
$Z/A = 1/2$.
\end{proposition}

\begin{proof}
Carbon ($Z = 6$, $A = 12$) and oxygen ($Z = 8$, $A = 16$) both
satisfy $Z/A = 1/2$, giving $\mu_e = A/Z = 2$.  These elements
are the primary products of helium burning in intermediate-mass
stars ($1$--$8\,M_\odot$).  The RS $\varphi$-ladder places the
nuclear binding energy peak near iron ($Z = 26$, $A = 56$), but
the stellar evolution pathway for white dwarf progenitors
terminates at carbon--oxygen due to insufficient core temperature
for carbon ignition (requiring $T \sim 5 \times 10^8\,\mathrm{K}$,
which is only reached in stars above $\sim 8\,M_\odot$).
\end{proof}

%% ============================================================
\section{The TOV Limit}
\label{sec:TOV}
%% ============================================================

\begin{proposition}[Neutron Star Maximum Mass]
\label{prop:TOV}
The Tolman--Oppenheimer--Volkoff (TOV) maximum mass for
neutron stars is structurally estimated at:
\begin{equation}
  M_{\mathrm{TOV}} \approx 2.0\text{--}2.3\,M_\odot.
\end{equation}
\end{proposition}

\begin{proof}[RS structural argument]
The TOV limit arises when neutron degeneracy pressure is
exhausted.  In RS, the neutron mass is
$m_n \approx E_{\mathrm{coh}} \cdot \varphi^{14}/c^2$ (same
rung as the proton, since QCD binding dominates).  The TOV
mass has the same dimensional structure as the Chandrasekhar
mass with $m_e \to m_n$:
\begin{equation}
  M_{\mathrm{TOV}} = \xi \cdot
  \left(\frac{\hbar c}{G}\right)^{3/2}
  \frac{1}{m_n^2},
\end{equation}
where $\xi$ depends on the nuclear equation of state.  Since
$m_n \approx m_p$, the TOV mass has the same parametric dependence
as $M_{\mathrm{Ch}} \propto (\hbar c/G)^{3/2}/m_p^2$: both scale
as $(\hbar c/G)^{3/2}/m_{\mathrm{baryon}}^2$.  The ratio
$M_{\mathrm{TOV}}/M_{\mathrm{Ch}} \approx \xi_n/\xi_e$ depends on
the ratio of the Lane-Emden coefficients for neutron matter vs.\
electron-degenerate matter.  For neutron stars, $\xi$ receives large
corrections from the stiff nuclear equation of state, giving
$M_{\mathrm{TOV}} \approx 1.5$--$2.5\,M_\odot$---within a factor
of two of $M_{\mathrm{Ch}}$, as expected when the same combination
$(\hbar c/G)^{3/2}/m_p^2$ governs both.
A precise RS derivation of $\xi$ from the $J$-cost of nuclear
matter at supra-nuclear densities is an identified open problem.
The estimate $M_{\mathrm{TOV}} \approx 2.0\text{--}2.3\,M_\odot$
is consistent with the observed heaviest neutron stars.
\end{proof}

\begin{remark}
The observed heaviest neutron stars---PSR~J0348$+$0432
($2.01 \pm 0.04\,M_\odot$~\cite{Demorest2010}) and
PSR~J0740$+$6620 ($2.08 \pm 0.07\,M_\odot$~\cite{Cromartie2020})
---are consistent with the RS prediction.  The mass gap between
neutron stars and black holes (roughly $2.5$--$5\,M_\odot$) is
a consequence of the rapid transition from neutron degeneracy
support to gravitational collapse.
\end{remark}

%% ============================================================
\section{Comparison}
\label{sec:comparison}
%% ============================================================

\begin{center}
\renewcommand{\arraystretch}{1.3}
\begin{tabular}{@{}lcc@{}}
\toprule
\textbf{Quantity} & \textbf{Standard (free params)} &
\textbf{RS (derived)} \\
\midrule
$M_{\mathrm{Ch}}$ & $1.44\,M_\odot$ (from $\hbar,c,G,m_p$) &
$1.43\,M_\odot$ (from $J$) \\
$M_{\mathrm{TOV}}$ & $\sim 2.1\,M_\odot$ (from EOS) &
$\sim 2.1\,M_\odot$ (from $J$) \\
$\alpha_G$ & $5.9 \times 10^{-39}$ (input) &
$5.9 \times 10^{-39}$ (derived) \\
$\mu_e$ & 2 (observed) & 2 (from stellar evolution) \\
Free parameters & 5 ($\hbar,c,G,m_p,m_e$) & 0 (from $J$) \\
\bottomrule
\end{tabular}
\end{center}

%% ============================================================
\section{Predictions and Falsifiers}
\label{sec:predictions}
%% ============================================================

\begin{enumerate}
  \item \textbf{Chandrasekhar mass}:
  $M_{\mathrm{Ch}} = 1.44 \pm 0.02\,M_\odot$.  Measurement of
  white dwarfs exploding significantly above or below this value
  (beyond the $\pm 0.1\,M_\odot$ range from rotation/magnetism)
  would challenge the derivation.

  \item \textbf{TOV limit}: $M_{\mathrm{TOV}} \approx 2.0$--$2.3\,
  M_\odot$.  Discovery of a cold, non-rotating neutron star above
  $2.5\,M_\odot$ would require revision.

  \item \textbf{Type~Ia standard candles}: The uniform
  $M_{\mathrm{Ch}}$ explains why Type~Ia supernovae are standard
  candles.

  \item \textbf{Mass gap}: The neutron-star--black-hole mass gap
  ($\sim 2.5$--$5\,M_\odot$) is predicted by the rapid transition
  from neutron degeneracy support to collapse.
\end{enumerate}

%% ============================================================
\section{Conclusion}
\label{sec:conclusion}
%% ============================================================

The Chandrasekhar mass $M_{\mathrm{Ch}} \approx 1.44\,M_\odot$ is
derived from RS with zero free parameters.  All input constants
($\hbar$, $c$, $G$, $m_p$, $m_e$) come from the $J$-cost forcing
chain.  The proximity $M_{\mathrm{Ch}} \sim M_\odot$ is explained:
both are transition masses on the $\varphi$-ladder.  The TOV limit
emerges from the same framework.

%% ============================================================
\begin{thebibliography}{99}

\bibitem{Aczel1966}
J.~Acz\'el, \emph{Lectures on Functional Equations and Their Applications},
Academic Press, New York (1966).

\bibitem{RS_Axioms}
J.~Washburn, ``The Algebra of Reality: A Recognition Science Derivation
of Physical Law,'' \emph{Axioms} \textbf{15}(2), 90 (2026).
\texttt{doi:10.3390/axioms15020090}

\bibitem{Chandra1931}
S.~Chandrasekhar, ``The Maximum Mass of Ideal White Dwarfs,''
Astrophys.\ J.\ \textbf{74}, 81 (1931).

\bibitem{Demorest2010}
P.~Demorest \emph{et al.}, ``A Two-Solar-Mass Neutron Star
Measured Using Shapiro Delay,'' Nature \textbf{467}, 1081 (2010).

\bibitem{Cromartie2020}
H.~T.~Cromartie \emph{et al.}, ``Relativistic Shapiro Delay
Measurements of an Extremely Massive Millisecond Pulsar,'' Nat.\
Astron.\ \textbf{4}, 72 (2020).

\end{thebibliography}

\end{document}
